\section{Problem Formulation}

\paragraph{Objects.}
A \emph{context} $x \in \mathcal{X}$ is a set of operating conditions known
before a policy is activated: demand, signal timing, guidance penetration,
information lag, disruption regime and alternative capacity. A \emph{policy}
$a \in \mathcal{A} = \{\mathrm{P1},\dots,\mathrm{P4}\}$ maps individual vehicles
to routes according to its own published objective. Running policy $a$ in
context $x$ produces an outcome vector; we write $C(x,a)$ for the scalar
decision cost, system mean journey time, and treat the other two criteria by
dominance (Section~\ref{sec:criteria}).

\paragraph{The two levels.}
The decision studied here is the map
\begin{equation}
\pi : \mathcal{X} \to \mathcal{A},
\qquad
x \;\xrightarrow{\ \pi\ }\; a \;\xrightarrow{\ \text{policy } a\ }\;
\text{routes} \;\longrightarrow\; \text{traffic outcome}.
\end{equation}
Level~1 is $\pi$, the object of study. Level~2 is the policy's own route
assignment, which is fixed and not learned. This is not route prediction, not
vehicle-level action learning, not signal control, and not joint
routing--signal optimisation.

\paragraph{Reference points.}
Two quantities from the algorithm-selection literature~\cite{bischl2016aslib}
bound what any $\pi$ can achieve. The \emph{single-best policy} (SBS) is the one
fixed policy with the lowest mean cost, chosen on training contexts only:
\begin{equation}
a_{\mathrm{SBS}} \;=\; \arg\min_{a \in \mathcal{A}}\
\mathbb{E}_{x \in \mathcal{X}_{\mathrm{train}}}\,[\,C(x,a)\,].
\end{equation}
It is what a deployed fixed-policy system achieves, and the baseline any
selector must beat to be worth building. The \emph{virtual-best policy} (VBS)
chooses per context using that context's realised outcomes,
$C_{\mathrm{VBS}}(x) = \min_{a} C(x,a)$. \emph{VBS is retrospective: it is not
an algorithm, cannot be deployed, and is used only as an upper reference.}

\paragraph{Regret and headroom.}
The regret of selecting $a$ in context $x$, and the mean regret of a selector,
are
\begin{equation}
R(x,a) \;=\; C(x,a) - \min_{b \in \mathcal{A}} C(x,b),
\qquad
\bar{R}(\pi) \;=\; \mathbb{E}_x\,[\,R(x,\pi(x))\,].
\end{equation}
The \emph{headroom} is the total decision value available to any selector:
\begin{equation}
H \;=\; \bar{R}(a_{\mathrm{SBS}})
   \;=\; \mathbb{E}_x\big[\,C(x,a_{\mathrm{SBS}}) - \min_b C(x,b)\,\big].
\label{eq:headroom}
\end{equation}
$H$ decides whether context-dependent selection is worth attempting at all, and
it is measurable without building a selector. We report it in seconds per
vehicle and, where a percentage is given, always as
$H / \mathbb{E}_x[C(x,a_{\mathrm{SBS}})]$ --- the aggregate decision headroom
relative to the single-best fixed policy, with that denominator stated
alongside. It is a share of an achievable decision cost, not a claim that any
vehicle's travel time improves by that fraction.

\paragraph{Three distinct questions.}
The paper keeps apart three properties that are often merged:
\begin{center}
\textbf{complementarity} $\;\neq\;$ \textbf{decision value} $\;\neq\;$
\textbf{deployability.}
\end{center}
Complementarity asks whether $\arg\min_a C(x,a)$ varies with $x$. Decision value
asks how large $H$ is. Deployability asks whether an estimate $\hat{\pi}$ built
from finite data retains that value outside the conditions it was fitted on. A
portfolio can be strongly complementary, carry little decision value, and still
produce a selector that is unreliable under shift. Distinguishing the three is
what the remainder of the paper does.

\section{Methodology}

\subsection{Simulation Environment and Network}

All experiments run in SUMO~\cite{lopez2018sumo}; \CO{} comes from its HBEFA3
implementation~\cite{krajzewicz2015emissions}. Policies act through the
simulator's online control interface, so every routing decision is taken inside
the running simulation from information available at that instant.

Figure~\ref{fig:network} shows the environment: one origin--destination pair
served by three structurally different corridors between a common diverge and
merge, crossed by four minor streets carrying independent demand through the
same signals, so a policy that loads a corridor also delays traffic it does not
route. Because every junction permits straight-through movements only, the
loopless path set is exactly these three --- the candidate catalogue is
complete, not a $k$-shortest approximation.

Two structural choices were fixed before any evaluation run, both identified by
capacity validation. The destination portal gives each approach a dedicated
exit lane, so an alternative corridor is not penalised by junction priority
rather than by its own capacity; and the diverge carries no turn-radius speed
ceiling, which in schematic node geometry would cap the bypass at a value
belonging to the drawing rather than to the road.

\paragraph{Measured capacity.} Each corridor was driven to saturation alone and
its discharge measured: \NumSatArterial~veh/h/lane on the arterial and
\NumSatNorth~veh/h/lane on the slower street, consistent to within
\NumSatSpread{} across green ratios, and \NumCapBypass~veh/h on the
uninterrupted bypass --- below queue-discharge saturation flow, as free-flow
headways require. The demand levels of Table~\ref{tab:factors} follow from these
measurements rather than from a textbook constant.

\subsection{Routing Policy Portfolio}

Table~\ref{tab:portfolio} gives each policy's objective, information
requirement, intended strength and expected failure mode. \emph{The failure
modes were specified before any evaluation run}, and each experimental factor
exists to make one of them reachable: P1 $\rightarrow$ capacity shortage on the
shortest corridor; P2 $\rightarrow$ information lag and herding of the guided
cohort; P3 $\rightarrow$ unevenly distributed alternative capacity;
P4 $\rightarrow$ travel-time variability under disruption.

All four read traffic state from one shared observation buffer. A policy acting
at time $t$ may read link travel times sampled no later than $t-\Delta$ over a
120\,s window, falling back to free-flow values before any admissible sample
exists. No policy sees the future, and none sees any realised outcome of the
run it is part of. Policy parameters were frozen before any evaluation run: a
travel-time admissibility tolerance $\varepsilon=\NumEpsFrozen$ and a 120\,s
assignment window for P3, and a dispersion weight $\lambda=1.0$ for P4.

Each policy was validated before use by twelve known-answer tests covering the
observation lag, the free-flow fallback and each policy's defining behaviour. A
two-path integration test inside the simulator, with equal path lengths and a
2:1 capacity ratio, confirms that P3 splits demand \NumTwoPathRatio{} against
that ratio while P1 commits the entire cohort to one path.

\subsection{Decision Criteria}
\label{sec:criteria}

Three criteria (Table~\ref{tab:criteria}) capture different burdens.
\textbf{C1, mean journey time}, is the travel burden per vehicle, including
insertion delay: a policy that oversaturates the network produces vehicles that
cannot enter it, and excluding their wait would credit the policy for the queue
it caused. \textbf{C2, \CO{}}, is the environmental outcome. \textbf{C3, total
stopped delay}, is the stationary burden --- time spent below 0.1\,m/s in
queues, which is distinct from time spent travelling slowly.

All three are measured over every vehicle in the network, guided or not, with
intended departure inside the measurement window, so a policy cannot be
rewarded for helping guided vehicles at the expense of background traffic.

Policies are compared by \textbf{Pareto dominance first}. The scalar cost used
for regret is C1, which requires no weights. Scalarisation appears only as
sensitivity, through three preference profiles declared before any evaluation
run and applied to within-context min--max normalised criteria. No monetary
value of time, value of reliability or carbon price is used.

\subsection{Mechanistic Rule}

A selector that fires on raw demand explains nothing, so each context is
expressed through eight dimensionless descriptors computable before activation:
network saturation $x_1$; shortest-corridor saturation
$x_2 = D/\mathrm{cap}_C$; alternative capacity share $x_3$; effective green
ratio $x_4$; penetration $x_5$; information staleness
$x_6 = \Delta/T^{\mathrm{ff}}_C$, the lag as a fraction of the free-flow trip
time it describes; incident regime $x_7$; and controlled-flow saturation
$x_8 = p\,D/\mathrm{cap}_C$, the share of the shortest corridor's capacity that
the guided cohort alone would consume. Capacities are the measured values, so
$x_1$, $x_2$ and $x_8$ are genuine degrees of saturation rather than scaled
demand. A route-overlap index is omitted: the three corridors share only access
and egress edges, so overlap is structurally constant here and carries no
information.

The mechanistic rule, denoted B1, is an axis-aligned region rule on
$x_1$--$x_8$ of depth at most two, fitted on training contexts only. Its
thresholds are reported as \emph{brackets between adjacent sampled levels},
since a grid cannot locate a boundary more finely than its own spacing.

\subsection{Context-Aware Selector}

The learned selector predicts \emph{advantage}, not identity:
\begin{equation}
\Delta_a(x) \;=\; C(x, a_{\mathrm{SBS}}) - C(x,a),
\qquad \hat{\pi}(x) = \arg\max_a \hat{\Delta}_a(x),
\label{eq:advantage}
\end{equation}
with $a_{\mathrm{SBS}}$ determined on training contexts only. A positive
$\Delta_a$ means $a$ beats the fixed policy a deployed system would otherwise
run, so the learned quantity is the one a decision-maker needs. The model is
gradient-boosted decision trees~\cite{friedman2001gbm,ke2017lightgbm}, one
regressor per policy, with hyperparameters fixed before any evaluation run and
never tuned on a test fold. With a few hundred context-level instances and full
feedback from every policy, higher capacity is neither needed nor defensible;
we use no graph neural network, deep network, reinforcement learning or
contextual bandit.

The baseline ladder is B0, the single best fixed policy; B1, the mechanistic
rule; B2, a depth-3 decision tree; B3, multinomial logistic regression; B4, the
advantage selector; and B5, its selective variant, with the hindsight benchmark
VBS as reference. \textbf{The primary comparison is B4 against B1}: B0 is a low
bar, and beating it shows only that the decision varies with context at all.

\subsection{Validation and Distribution-Shift Framework}

Evaluation is leave-one-context-out. Inside each fold, standardisation, the
identity of the single best policy, tree thresholds, conformal quantiles and
the support envelope are recomputed from training data alone. No seed of a
held-out context appears in training, and no realised outcome is used as a
feature --- the eight descriptors are functions of the context definition and
cannot be computed from an outcome.

\paragraph{Two gates.} B5 acts only if both pass. The \emph{support gate}
measures mean distance to the five nearest training contexts in standardised
descriptor space and abstains above the 95th percentile of the training fold's
own such distances; it is a feature-space extrapolation detector. The
\emph{confidence gate} is a split-conformal interval on the predicted advantage
($\alpha=0.10$, calibrated on 30\% of the training fold); the selector acts only
when the lower bound on the advantage over the fixed policy is positive, and it
is calibrated under exchangeability, hence in-distribution. If either gate
fails, the fixed policy is used.

Neither gate detects concept shift, and the conformal guarantee does not survive
covariate shift~\cite{shimodaira2000,tibshirani2019covshift}. We therefore
report gate behaviour rather than coverage guarantees outside the training
distribution, and the support gate is not presented as a general
out-of-distribution detector.

\subsection{Statistical Definitions and Evaluation Metrics}
\label{sec:stats}

The experimental unit is the context, not the vehicle. An advantage of policy
$a$ over policy $b$ in a context is \emph{resolved} only when the paired seed
differences $d_i = C_i(x,a) - C_i(x,b)$ satisfy
\begin{equation}
|\overline{d}| \;>\; 2\,\mathrm{SE}(\overline{d}),
\qquad
\mathrm{SE}(\overline{d}) = s_d / \sqrt{n},
\label{eq:resolved}
\end{equation}
with $n$ the number of paired seeds. Differences that fail this criterion are
treated as ties throughout, so no claim rests on a difference the seeds cannot
separate. We report this resolution criterion rather than $p$-values: it is a
descriptive threshold on paired replications, and we make no claim of
statistical significance in the hypothesis-testing sense.

Selector quality is reported as decision regret (mean, median, maximum and
CVaR$_{90}$ over contexts), the share of contexts within the seed noise floor,
the share exactly optimal, and the fraction of the available headroom closed.
Selection accuracy --- how often the selector names the true best policy --- is
reported separately and never used as a proxy for regret.
